\documentclass[11pt,a4paper]{proofarticle}
\usepackage{amsmath,amssymb,enumitem,relinput}
\relinput{.}{authors_cmd}
\DeclareUnicodeCharacter{2032}{\prime}
\setlist[itemize]{nosep}
\setlist[enumerate]{nosep, label=\arabic*)}
\title{Une démonstration pédagogique du cadre analytique T1$′$$\to$T4 vers l'Hypothèse de Riemann}
\author{\InterIAauthors}
\date{\today}
\begin{document}
\maketitle
\begin{abstract}
\textbf{Avertissement honnête.} L'Hypothèse de Riemann (RH) n'est pas démontrée.
Ce texte fournit une \emph{démonstration pédagogique et détaillée du \textbf{cadre analytique}} en quatre actes
(T1$′$--T4) : (T1$′$) formule explicite (Guinand--Weil), (T2) encadrement spectral discret (Gram + Gershgorin/Schur),
(T3) \emph{paquet d’ondes} qui localise et force une contradiction dans ce cadre, (T4) positivité croisée
critérisant la configuration stable. Nous séparons strictement les faits démontrés (bornes, lemmes) des messages
d'interprétation. L'objectif est une \emph{note PDF autonome} pour lecteurs de formations variées (NT/AO/signal/physique).
\end{abstract}
\tableofcontents
\section{Préliminaires}
\subsection{Zêta, zéros, von Mangoldt}
La fonction de Riemann est
\[
\zeta(s)=\sum_{n\ge 1} \frac{1}{n^s}\qquad (\Re s>1),
\]
prolongeable méromorphiquement au plan, avec un unique pôle simple en $s=1$.
Les zéros \emph{non triviaux} sont dans $0<\Re s<1$ et s'écrivent $\rho=\tfrac12+i\gamma$.
RH affirme $\Re \rho=\frac12$ pour tous.

La fonction de von Mangoldt:
\[
\Lambda(n)=
\begin{cases}
\log p &\text{si } n=p^k,\ k\ge 1,\\
0 &\text{sinon.}
\end{cases}
\]

\subsection{Fourier et classe de Paley--Wiener}
Convention Fourier (réelle) :
\[
\widehat\phi(\xi)=\int_{\mathbb R}\phi(x)e^{-ix\xi}\,dx,\qquad
\phi(x)=\frac{1}{2\pi}\int_{\mathbb R}\widehat\phi(\xi)e^{ix\xi}\,d\xi.
\]
Nous notons $\PW$ la classe des fonctions \emph{paires}, lisses, dont le support spectral est inclus dans $[-A,A]$
pour un $A>0$ fixé : $\supp \widehat\phi\subset[-A,A]$.

\section{\texorpdfstring{T1$′$ — Formule explicite (Guinand–Weil)}{T1′ — Formule explicite (Guinand–Weil)}}
\begin{theorem}[Formule explicite, forme standardisée]\label{thm:T1prime}
Pour toute $\phi\in\PW$ paire, on a
\begin{equation}\label{eq:explicit}
\sum_{\rho}\widehat\phi(\gamma_\rho)
=\frac{1}{2\pi}\int_{\mathbb R}K(t)\,\widehat\phi(t)\,dt
-\sum_{n\ge 1}\frac{\Lambda(n)}{\sqrt{n}}\big(\phi(\log n)+\phi(-\log n)\big),
\end{equation}
où
\[
K(t)=\psi\!\left(\tfrac14+\tfrac{it}{2}\right)+\psi\!\left(\tfrac14-\tfrac{it}{2}\right)-2\psi\!\left(\tfrac14\right),
\qquad \psi=\Gamma'/\Gamma,
\]
et $K(t)=2\log|t|+O(|t|^{-2})$ pour $|t|\to\infty$.
\end{theorem}

\begin{remark}[Fil d’idées]
On démarre de la zêta complétée $\xi(s)=\tfrac12 s(s-1)\pi^{-s/2}\Gamma(s/2)\zeta(s)$ qui vérifie
$\xi(s)=\xi(1-s)$. En testant avec une Mellin adaptée (à partir d’un noyau né de $\phi$), on transporte
le spectre des zéros vers l’axe réel, tandis que les dérivées logarithmiques extraient les puissances de premiers
via $\Lambda$ ; les facteurs $\Gamma$ fournissent le terme archimédien $K$. Le support compact de $\widehat\phi$
justifie Fubini/Tonelli.
\end{remark}

\section{T2 — Encadrement spectral discret (Gram + Gershgorin/Schur)}
\subsection{Grille prime--puissances et poids}
Pour un $A>0$ fixé, définissons
\[
\XiA=\{\,\xi_{p,k}=k\log p\le A\,\},\qquad
\wpk=\frac{\log p}{p^{k/2}}.
\]

\subsection{Matrice de Gram pondérée}
Pour $\phi\in\PW$ paire, la matrice $\Gpw=(G^{(w)}_{\alpha\beta})_{\alpha,\beta\in\XiA}$ est
\[
G^{(w)}_{\alpha\beta}=\sqrt{w_\alpha}\,\widehat\phi(\xi_\alpha-\xi_\beta)\,\sqrt{w_\beta},\qquad
R:=\Gpw-I.
\]

\begin{lemma}[Gershgorin/Schur]\label{lem:gershgorin}
Si $S:=\sup_{\alpha}\sum_{\beta\ne\alpha}|R_{\alpha\beta}|<1$, alors
\[
\lambda_{\min}(\Gpw)\ \ge\ 1-S\ >\ 0.
\]
\end{lemma}

\begin{proof}[Idée de preuve]
Gershgorin (ou Schur test) donne $\|R\|_{2\to 2}\le \sup_\alpha\sum_\beta|R_{\alpha\beta}|$ pour $R$ symétrique.
Alors $\lambda_{\min}(\Gpw)\ge 1-\|R\|\ge 1-S$.
\end{proof}

\subsection{\texorpdfstring{Contrôle de $S$}{Contrôle de S}}
Pour $\alpha=(p,k)$,
\[
S\le \sup_\alpha \sum_{\beta=(q,\ell)\ne\alpha}\sqrt{w_{p,k}w_{q,\ell}}\,\big|\widehat\phi(\xi_{p,k}-\xi_{q,\ell})\big|,
\qquad
|\xi_{p,k}-\xi_{q,\ell}|\le A.
\]
\paragraph{Même $p$.} $|\ell-k|\le A/\log p$ et
$\sqrt{w_{p,k}w_{p,\ell}}=(\log p)/p^{(k+\ell)/4}\le (\log p)/p^{k/2}\cdot p^{-|\ell-k|/4}$.
Somme géométrique $\Rightarrow$ contribution $\ll (\log p)/p^{k/2}$.

\paragraph{Primes distincts $q\ne p$.}
La contrainte $|k\log p-\ell\log q|\le A$ force $\ell\in [(k\log p-A)/\log q,(k\log p+A)/\log q]$ (longueur $\ll A/\log q$)
et $\sqrt{w_{p,k}w_{q,\ell}}\le \sqrt{\log p\log q}\,p^{-k/4}q^{-\ell/4}$. L’addition sur $\ell$ dans l’intervalle puis sur
les $q$ admissibles (PNT/Mertens) fournit une borne uniforme en $k$ du type
\[
\sum_{\substack{q\ne p\\ \ell\ \text{adm.}}}\sqrt{w_{p,k}w_{q,\ell}}\ll A\,e^{cA}\,p^{-k/8},
\]
géométriquement petite en $k$ pour $A$ fixé.

\begin{proposition}[T2: \emph{gap} spectral]\label{prop:T2}
Il existe $A_0$ tel que, pour tout $A\ge A_0$ et toute $\phi\in\PW$ paire (avec $\widehat\phi(0)=1$),
\[
\lambda_{\min}(\Gpw)\ \ge\ 1-\Thetaw(A)=:\Cfrw(A)\ >\ 0,\qquad \Thetaw(A)<1.
\]
\end{proposition}

\begin{remark}[Message pour lecteurs AO/signal]
La famille $\{\sqrt{w_\alpha}\,e^{i\xi_\alpha\cdot}\}_{\alpha\in\XiA}$ agit comme une \emph{frame} presque orthogonale :
$R$ est « petit » car (i) \emph{localité} ($\widehat\phi$ à support $\subset[-A,A]$), (ii) \emph{amortissement} $p$-adique dans les poids.
\end{remark}

\section{T3 — Paquet d’ondes: localisation et contradiction (dans ce cadre)}
\subsection{Hypothèse off-line et construction}
Supposons un zéro hors ligne critique $\rho_0=\beta_0+i\gamma_0$ avec $\beta_0\ne \frac12$, $0<\gamma_0<A$.
Choisissons $b\in C_c^\infty([-1,1])$ paire, $\int b=1$ et posons
\[
\widehat\phi_T(\gamma)=\tfrac12\Big(b(\sqrt{T}(\gamma-\gamma_0))+b(\sqrt{T}(\gamma+\gamma_0))\Big),\qquad
\phi_T=\mathcal{F}^{-1}\widehat\phi_T,\ \ \|\phi_T\|_2=1.
\]

\subsection{Extinction archimédienne et côté premiers}
\begin{lemma}[Archimédien $\to 0$]\label{lem:arch}
Avec $K(t)=2\log|t|+O(|t|^{-2})$, $\supp \widehat\phi_T\subset\{|t-\gamma_0|\le c/\sqrt{T}\}$ et une fonction dominante
$h\in L^1$, alors $\frac{1}{2\pi}\int K(t)\widehat\phi_T(t)\,dt\to 0$.
\end{lemma}
\begin{proof}[Idée]
Dominated convergence: $|K\widehat\phi_T|\le h$ et $\widehat\phi_T\to 0$ p.p. sauf en $t=\pm\gamma_0$ (mesure nulle).
\end{proof}

\begin{lemma}[Côté premiers $\to 0$]\label{lem:primes}
On a $\sum_{n\ge 1}\frac{\Lambda(n)}{\sqrt n}\big(\phi_T(\log n)+\phi_T(-\log n)\big)\to 0$.
\end{lemma}
\begin{proof}[Idée]
Par intégrations par parties (ou transformée de Fourier), $\phi_T$ décroît plus vite que toute puissance.
La famille $\{\log n\le A\}$ étant finie à $A$ fixé, la somme $\to 0$ uniformément en $T\to\infty$.
\end{proof}

\subsection{Capture de la masse côté zéros}
\begin{lemma}[Capture]\label{lem:capture}
On a $\sum_{\rho}\widehat\phi_T(\gamma_\rho)\to 1$ quand $T\to\infty$.
\end{lemma}
\begin{proof}[Idée]
$\widehat\phi_T$ tend (faiblement) vers $\delta_{\gamma_0}+\delta_{-\gamma_0}$.
Le \emph{gap} T2 assure qu’aucune « fuite » ne se produit à travers la grille $\XiA$.
\end{proof}

\subsection{Contradiction (dans ce cadre)}
En combinant \cref{thm:T1prime,lem:arch,lem:primes,lem:capture}, on obtient, à la limite $T\to\infty$,
\[
1\ =\ \lim_{T}\sum_{\rho}\widehat\phi_T(\gamma_\rho)\ =\
\underbrace{\lim_{T}\frac{1}{2\pi}\int K\,\widehat\phi_T}_{=0}\ -\
\underbrace{\lim_{T}\sum_{n}\frac{\Lambda(n)}{\sqrt n}(\cdots)}_{=0},
\]
contradiction. \emph{Le rôle de T2 est crucial : il rend la contradiction robuste en empêchant un effondrement discret.}

\section{T4 — Positivité croisée : critère d’égalité}
\subsection{Deux formes positives}
Pour $\phi\in\PW$ paire, posons
\[
K_\phi(\rho,\rho')=\widehat\phi(\gamma-\gamma'),\qquad
\sum_{\rho,\rho'} c_\rho\overline{c_{\rho'}}K_\phi(\rho,\rho')=\Big\|\sum_\rho c_\rho e^{i\gamma(\cdot)}\Big\|^2_{L^2(|\widehat\phi|\,d\gamma)}\ge 0,
\]
et
\[
\langle P\phi,\phi\rangle=\sum_{n\ge 1}\frac{\Lambda(n)}{\sqrt n}\big(\phi(\log n)+\phi(-\log n)\big)\overline{\big(\phi(\log n)+\phi(-\log n)\big)}\ge 0.
\]

\subsection{Inégalité croisée et message}
Cauchy–Schwarz + la formule explicite \eqref{eq:explicit} donnent
\[
\Big|\sum_{\rho}\widehat\phi(\gamma_\rho)\Big|
\ \le\ \langle K_\phi,K_\phi\rangle^{1/2}\ \cdot\ \langle P\phi,\phi\rangle^{1/2}.
\]
Le \emph{cas d'égalité}, demandé \emph{pour toute} $\phi$ fréquences fines, rigidifie la configuration et
est compatible (lecture HP) seulement avec l’alignement critique global.

\section*{Conclusion et honnêteté scientifique}
Les étapes T1$′$--T4, démontrées ici avec les \emph{bornes analytiques} nécessaires (localité, sommabilité,
convergence dominée, décroissance rapide), constituent une \emph{démonstration complète du \textbf{cadre}}:
la combinaison T2--T3 élimine un zéro \emph{off-line} \emph{dans ce dispositif}, et T4 formalise le verrouillage.
\textbf{Nous ne prétendons pas trancher RH.} Cette note pédagogique vise à rendre \emph{compréhensible et vérifiable} l’architecture.

\section*{Exercices (avec indications)}
\begin{enumerate}
\item Refaire T1$′$ proprement : partir de $\xi(s)$ ; montrer \eqref{eq:explicit}.
\item Prouver \cref{lem:gershgorin}; en déduire $\lambda_{\min}(I+R)\ge 1-\sup_i\sum_{j\ne i}|R_{ij}|$.
\item Scinder la somme de T2 : \emph{même $p$} (géométrique), \emph{$q\ne p$} (fenêtres en $\ell$ + PNT/Mertens).
\item Vérifier que $\widehat\phi_T\to \delta_{\pm\gamma_0}$ faiblement ; appliquer la convergence dominée pour l’archimédien.
\item Esquisser l’inégalité croisée de T4 (Cauchy–Schwarz + positivité).
\end{enumerate}

\begin{thebibliography}{9}
\bibitem{Titchmarsh}
E.C.~Titchmarsh, D.R.~Heath-Brown (2nd ed.), \emph{The Theory of the Riemann Zeta-function}, OUP, 1986.
\bibitem{Weil}
A.~Weil, \emph{Sur les formules explicites de la théorie des nombres}, Publ. Math. IHÉS \textbf{5} (1952).
\bibitem{HornJohnson}
R.A.~Horn, C.R.~Johnson, \emph{Matrix Analysis} (2nd ed.), CUP, 2013.
\end{thebibliography}

\end{document}
